\chapter{Videonystagmography (VNG)}
\section*{What Is This Test?} VNG records eye movements to assess the vestibular system and causes of dizziness or vertigo.\section*{Alternative Names} VNG; videonystagmography; vestibular eye-movement testing.\section*{Principle} Infrared goggles track nystagmus during visual, positional, and caloric maneuvers.\section*{Why This Test Is Ordered} It evaluates vertigo, imbalance, suspected inner-ear disease, and abnormal eye movements.\section*{Primary vs Secondary Test} History and examination are primary; VNG is a selected secondary vestibular test.\section*{How This Test Is Performed} The patient follows targets, changes position, and may receive warm or cool air or water in the ear canal while eye movements are recorded.\section*{What Preparation Is Needed} Avoid alcohol and follow instructions about caffeine and medicines; do not stop medicines without guidance. Arrange help home if dizziness is expected.\section*{Understanding Results} Patterns may support peripheral or central vestibular dysfunction but require clinical correlation.\section*{Follow-up Tests} Audiometry, imaging, neurologic assessment, or vestibular rehabilitation may follow.\section*{References} \begin{enumerate}\item MedlinePlus. Videonystagmography (VNG).\item American Academy of Audiology. Vestibular assessment guidance.\end{enumerate}\clearpage\section*{Understanding Results and Follow-up}
The laboratory report should identify the specimen, method, units, reference interval or decision limit, and whether the result is positive, negative, abnormal, or inconclusive. Reference intervals vary with age, sex, pregnancy, specimen type, assay, and laboratory; a result outside a range is not automatically a diagnosis, and a result within a range does not always exclude disease. Discuss unexpected results with the ordering clinician, who can decide whether repeat or confirmatory testing, treatment, or referral is appropriate. Do not change medicines or delay urgent care based on an isolated result.
\section*{Regulatory Status and Authorization Date}This chapter describes a test category, not a specific commercial product. FDA, CDSCO, EU IVDR, and MHRA approval or registration dates therefore cannot be assigned without the assay manufacturer, product name, instrument, and intended use. For laboratory-developed tests, an individual agency approval date may not exist. See the Preface for the interpretation of regulatory status.
\clearpage
